\documentclass[11pt]{article}
\usepackage[a4paper,margin=27mm]{geometry}
\usepackage{amsmath,amssymb,mathtools}
\usepackage{fontspec,unicode-math}
\setmainfont{Latin Modern Roman}
\setmathfont{Latin Modern Math}
\usepackage{xurl}
\usepackage[hidelinks]{hyperref}
\newcommand{\sref}[2]{\hyperlink{source:#1:#2}{[S#2]}}
\newcommand{\cataloguescope}[1]{\paragraph{Mathematical question.}#1}
\hypersetup{pdftitle={Problem statement}}
\begin{document}
\section{Carmichael's totient function conjecture}\label{statement}

\subsection{Definitions and mathematical statement}

\paragraph{Shared notation.}$\mathbb N=\{1,2,\ldots\}$, and $\mathbb Z,\mathbb Q,\mathbb R,\mathbb C$ denote the integers, rationals, reals and complex numbers. Any other conventions are specified locally.


Euler's totient is $\varphi(n)=\#\{a\in\mathbb N:1\le a\le n,\ \gcd(a,n)=1\}$ for $n\in\mathbb N$. For $m\in\mathbb N$, its multiplicity is $A(m)=\#\{n\in\mathbb N:\varphi(n)=m\}$.
\paragraph{Statement recorded in the supplied source.}
Is $A(m)\ge2$ for every $m\in\varphi(\mathbb N)$? Equivalently, for every $n\in\mathbb N$, does there exist $n'\in\mathbb N$, $n'\ne n$, with $\varphi(n')=\varphi(n)$? \sref{statement}{2}
\subsection{Short English statement}
Does every totient value come from at least two distinct positive integers?
\subsection{Sources}
\begin{description}
\item[\hypertarget{source:statement:1}{S1}] ULAM.AI and the UnsolvedMath Contributors, UnsolvedMath: A Curated Collection of Open Mathematics Problems, catalogue entry. CC BY 4.0. \url{https://huggingface.co/datasets/ulamai/UnsolvedMath}.
\item[\hypertarget{source:statement:2}{S2}] Mithun Kumar Das, Pramod Eyyunni and Bhuwanesh Rao Patil, \emph{Sparse subsets of the natural numbers and Euler\textquotesingle s totient function}, arXiv:1907.09847v2 (2020) Introduction, p. 3, statement before Theorem 1.6. \url{https://arxiv.org/pdf/1907.09847}.
\item[\hypertarget{source:statement:3}{S3}] Wikipedia contributors, \emph{List of unsolved problems in mathematics}, revision 1366636767 (2026) Number theory, Carmichael's totient function conjecture. \url{https://en.wikipedia.org/w/index.php?title=List_of_unsolved_problems_in_mathematics&oldid=1366636767#Number_theory}.
\end{description}
\label{end:statement}
\end{document}
